\chapter{SAMPLE SOURCE CODE IMPLEMENTATION}
\label{app:code}

This Appendix provides clean, representative excerpts of the core algorithmic modules developed for the Wasserstein-guided multi-task distillation framework.

\section{Sinkhorn-Wasserstein Optimal Transport Solver}
The following module implements the vectorized, GPU-accelerated entropic Optimal Transport loss with cosine ground cost computation:

\begin{singlespace}
\begin{verbatim}
import torch
import torch.nn as nn
import torch.nn.functional as F

class SinkhornWassersteinLoss(nn.Module):
    """
    Computes cross-vocabulary entropic Sinkhorn-Wasserstein distance
    between student acoustic representations and teacher token embeddings.
    """
    def __init__(self, epsilon=0.10, max_iter=10):
        super(SinkhornWassersteinLoss, self).__init__()
        self.epsilon = epsilon
        self.max_iter = max_iter

    def forward(self, student_feats, teacher_feats, student_mask=None):
        # Normalize representations to unit hypersphere
        s_norm = F.normalize(student_feats, p=2, dim=-1) # (B, N, D)
        t_norm = F.normalize(teacher_feats, p=2, dim=-1) # (B, M, D)

        # Pairwise Cosine Ground Cost Matrix C_ij = 1 - cos(s_i, t_j)
        cost_matrix = 1.0 - torch.bmm(s_norm, t_norm.transpose(1, 2)) # (B, N, M)
        cost_matrix = torch.clamp(cost_matrix, min=0.0, max=2.0)

        B, N, M = cost_matrix.shape

        # Uniform empirical target measures
        r = torch.ones(B, N, device=student_feats.device) / N
        c = torch.ones(B, M, device=student_feats.device) / M

        # Exponential Gibbs Kernel
        kernel = torch.exp(-cost_matrix / self.epsilon) + 1e-9

        # Iterative Sinkhorn-Knopp Matrix Scaling
        b = torch.ones(B, M, device=student_feats.device) / M
        for _ in range(self.max_iter):
            # Row normalization
            a = r / (torch.bmm(kernel, b.unsqueeze(-1)).squeeze(-1) + 1e-9)
            # Column normalization
            b = c / (torch.bmm(kernel.transpose(1, 2), a.unsqueeze(-1)).squeeze(-1) + 1e-9)

        # Optimal Transport Coupling Plan P*
        P_star = a.unsqueeze(-1) * kernel * b.unsqueeze(1) # (B, N, M)

        # Frobenius Inner Product Loss: <P*, C>
        ot_loss = torch.sum(P_star * cost_matrix, dim=(-2, -1)).mean()
        return ot_loss
\end{verbatim}
\end{singlespace}

\section{Dual-Head Multi-Task Conformer Architecture}
The following class defines the multi-task prediction heads and composite loss computation:

\begin{singlespace}
\begin{verbatim}
class MultiTaskConformerStudent(nn.Module):
    def __init__(self, vocab_size=1024, num_emotions=4, hidden_dim=256):
        super().__init__()
        self.encoder = ConformerBackbone(input_dim=80, d_model=hidden_dim, 
                                         num_layers=6, conv_kernel_size=15)
        self.linear_proj = nn.Linear(hidden_dim, 2048) # Teacher manifold projection

        # ASR Head
        self.asr_head = nn.Linear(hidden_dim, vocab_size)

        # SER Affective Head
        self.ser_head = nn.Sequential(
            nn.Linear(hidden_dim, 128),
            nn.GELU(),
            nn.Dropout(0.2),
            nn.Linear(128, num_emotions)
        )

    def forward(self, mel_spectrogram, lengths=None):
        # 4x subsampled Conformer features: (B, T', 256)
        latents, pad_mask = self.encoder(mel_spectrogram, lengths)

        # Acoustic transcription logits
        asr_logits = self.asr_head(latents)

        # Temporal mean pooling over unmasked frames for SER
        if pad_mask is not None:
            mask_expanded = (~pad_mask).unsqueeze(-1).float()
            pooled = (latents * mask_expanded).sum(dim=1) / mask_expanded.sum(dim=1).clamp(min=1)
        else:
            pooled = latents.mean(dim=1)

        emotion_logits = self.ser_head(pooled)
        projected_latents = self.linear_proj(latents)

        return asr_logits, emotion_logits, projected_latents
\end{verbatim}
\end{singlespace}
